\section{信頼性と故障}

\subsection{平成23年度 問13}
\label{q:h23-reg-apt-13}

\noindent\textbf{原典：} sources/H23/aptitude.pdf、PDF物理14頁。

\textbf{公式正答：} 3（single）。

\par\textit{原典本文の校正と独自解説は未作成。分類のみの登録項目。}

\section{公益・説明責任}

\subsection{平成23年度 問15}
\label{q:h23-reg-apt-15}

\noindent\textbf{原典：} sources/H23/aptitude.pdf、PDF物理16頁。

\textbf{公式正答：} 1（single）。

\paragraph{問題の要旨} 社会的合意を形成する際の進行役・専門家・第三者の役割から、不適切な説明を選ぶ。

\Needspace{4\baselineskip}\paragraph{確認した条件・記述（独自要約）}\mbox{}\par

\noindent\textbf{方向：} 最も不適切なもの。\par

\noindent\textbf{状況：} 利害・価値観の異なる関係者間で社会的合意を形成する。\par

\paragraph{問われる核心} 合意形成における専門助言と中立な手続支援

\paragraph{必要知識} 関係者分析、利害、進行役、第三者の中立性、技術的判断と価値判断。

\paragraph{考え方・正答を選ぶ理由（査読待ち初稿）} 合意形成は専門家が一方的に正解を説明する行為ではなく、情報・利害・価値・手続を調整する過程である。担当者の専門性と進行役の中立性・能力を同一視した1が不適切である。第三者を置けば必ず全員一致するという意味でもない。

\Needspace{4\baselineskip}\paragraph{選択肢の検討（査読待ち初稿）}\mbox{}\par

\noindent\textbf{1：設問条件に一致} 業務担当技術者自身が、あらゆる問題に対応する進行役となる。。担当者は事業への利害も持つ。専門知識だけで中立な進行能力やあらゆる対立の処理を保証できない。\par

\noindent\textbf{2：設問条件に不一致} 技術的役割と会議進行・意見調整の役割を分けて検討する。。科学技術の助言と公正な発言機会の確保は異なる能力・役割である。\par

\noindent\textbf{3：設問条件に不一致} 関係者とその利害を把握する。。表面の賛否だけでなく、守りたい利益や懸念を把握して調整する。\par

\noindent\textbf{4：設問条件に不一致} 中立な第三者の支援を得る。。中立性を担保した第三者は手続の公正さを支援し得る。\par

\noindent\textbf{5：設問条件に不一致} 進行支援の必要性を事前評価する。。関係者・対立点・実施可能性を調べて適切な支援を設計する。\par

\paragraph{再利用する考え方} 合意形成における専門助言と中立な手続支援について、対象・要件・例外・時点を分け、類似用語や無条件の断定を避ける。

\paragraph{共通理論} \hyperref[theory:compliance:consensus]{第\ref*{theory:compliance:consensus}節}

\par 参考：\cite{mlit_mediation2006,mlit_pi_thirdparty}。

\par 難易度 3（校準初稿）、優先度未評価。

\subsection{平成28年度 問6}
\label{q:h28-reg-apt-06}

\noindent\textbf{原典：} sources/H28/aptitude.pdf、PDF物理6頁。

\textbf{公式正答：} 1（single）。

\paragraph{問題の要旨} 公衆への説明責任について4記述の正誤の組を選ぶ。

\Needspace{4\baselineskip}\paragraph{確認した条件・記述（独自要約）}\mbox{}\par

\noindent\textbf{ア：} 十分な説明に基づく理解・信頼が重要。\par

\noindent\textbf{イ：} 専門知識を持たない公衆に分かるよう専門家が説明する必要がある。\par

\noindent\textbf{ウ：} 説明責任と守秘義務の双方が必要。\par

\noindent\textbf{エ：} 知る権利に応答する説明を、自分の意見の押付けと区別する。\par

\noindent\textbf{並び：} ア／イ／ウ／エ\par

\paragraph{問われる核心} 専門知識の非対称性と理解可能な説明

\paragraph{必要知識} 説明責任、信頼、知る権利、秘密保持、同意と意見の押付け。

\paragraph{考え方・正答を選ぶ理由（査読待ち初稿）} 4項とも適切であり1である。イは専門外の人にも理解できる説明を行う責任で、同じ専門能力を要求したり、公衆の価値判断を無視したりする趣旨ではない。説明には理解・質問・検討の機会が伴う。守秘を無条件の隠蔽へ、説明を無制限の開示へ変えない。

\Needspace{4\baselineskip}\paragraph{選択肢の検討（査読待ち初稿）}\mbox{}\par

\noindent\textbf{1：設問条件に一致} 正／正／正／正。全項の判断に一致する。\par

\noindent\textbf{2：設問条件に不一致} 誤／誤／正／正。ア・イの正誤が違う。各記述の定義・条件を先に判定してから表に対応させる。\par

\noindent\textbf{3：設問条件に不一致} 正／正／誤／正。ウの正誤が違う。各記述の定義・条件を先に判定してから表に対応させる。\par

\noindent\textbf{4：設問条件に不一致} 正／正／誤／誤。ウ・エの正誤が違う。各記述の定義・条件を先に判定してから表に対応させる。\par

\noindent\textbf{5：設問条件に不一致} 正／誤／誤／正。イ・ウの正誤が違う。各記述の定義・条件を先に判定してから表に対応させる。\par

\paragraph{再利用する考え方} 専門知識の非対称性と理解可能な説明について、対象・要件・例外・時点を分け、類似用語や無条件の断定を避ける。

\paragraph{共通理論} \hyperref[theory:compliance:accountability]{第\ref*{theory:compliance:accountability}節}

\par 参考：\cite{ipej_ethics2011,ieej_conduct2007}。

\par 難易度 3（校準初稿）、優先度未評価。

\section{知的財産}

\subsection{平成27年度 問12}
\label{q:h27-reg-apt-12}

\noindent\textbf{原典：} sources/H27/aptitude.pdf、PDF物理12頁。

\textbf{公式正答：} 2（single）。

\paragraph{問題の要旨} 学術論文等の引用について適切な説明を選ぶ。

\Needspace{4\baselineskip}\paragraph{確認した条件・記述（独自要約）}\mbox{}\par

\noindent\textbf{時点：} 平成27年。\par

\noindent\textbf{方向：} 最も適切なもの。\par

\paragraph{問われる核心} 引用という要件付きの権利制限

\paragraph{必要知識} 著作権法32・43・48条、公表、公正な慣行、正当な範囲、出所、翻訳と翻案。

\paragraph{考え方・正答を選ぶ理由（査読待ち初稿）} 引用は権利者の許諾を必要とする利用一般に対する要件付きの例外であり2が適切である。公表されたから自由という説明、引用だから量・目的を問わないという説明はともに誤る。訳文の利用も引用要件を満たす必要がある。

\Needspace{4\baselineskip}\paragraph{選択肢の検討（査読待ち初稿）}\mbox{}\par

\noindent\textbf{1：設問条件に不一致} 承諾のない引用は全て著作権侵害である。。著作権法32条の要件を満たす引用という権利制限がある。\par

\noindent\textbf{2：設問条件に一致} 適法な引用の条件を満たせば、承諾なしに引用できる。。公表、公正な慣行、目的上正当な範囲などを満たす場合が対象で、出所も明示する。\par

\noindent\textbf{3：設問条件に不一致} 公表済みの著作物は自由に複製できる。。公表と権利放棄は別であり、一般に無制限の複製を許さない。\par

\noindent\textbf{4：設問条件に不一致} 引用は学術研究目的だけに限られる。。報道・批評・研究その他の引用目的で、研究だけとの限定は誤り。\par

\noindent\textbf{5：設問条件に不一致} 外国語の著作物は原語のままでしか引用できない。。適法引用の範囲では43条による翻訳利用が可能であり、原語に限定されない。\par

\paragraph{再利用する考え方} 引用という要件付きの権利制限について、対象・要件・例外・時点を分け、類似用語や無条件の断定を避ける。

\paragraph{共通理論} \hyperref[theory:compliance:citation]{第\ref*{theory:compliance:citation}節}

\par 参考：\cite{bunka_citation_faq,kyushu_copyright2015}。

\paragraph{関連問題} \hyperref[q:r06-reg-apt-03]{R06-REG-APT-03}

\par 難易度 3（校準初稿）、優先度未評価。

\subsection{平成28年度 問11}
\label{q:h28-reg-apt-11}

\noindent\textbf{原典：} sources/H28/aptitude.pdf、PDF物理11頁。

\textbf{公式正答：} 4（single）。

\paragraph{問題の要旨} 知的財産権・特許の発明・実用新案の考案・ソフトウェアの4説明の組を選ぶ。

\Needspace{4\baselineskip}\paragraph{確認した条件・記述（独自要約）}\mbox{}\par

\noindent\textbf{ア：} 法令に定める権利や法律上保護される利益に係る権利を含む知的財産権の定義。\par

\noindent\textbf{イ：} 発明は自然法則を利用した技術的思想の創作のうち高度のもの。\par

\noindent\textbf{ウ：} 考案も自然法則を利用した技術的思想の創作で、高度を要求しない。\par

\noindent\textbf{エ：} 抽象的な人為的取決め・暗号作成・計算は発明に当たらず、プログラムでも自然法則を利用する必要がある。単なるプログラムリストの提示は発明に当たらない。\par

\noindent\textbf{並び：} ア／イ／ウ／エ。平成28年。\par

\paragraph{問われる核心} 技術的創作と単なる規則・情報の区別

\paragraph{必要知識} 知的財産基本法2条、特許法2条、実用新案法2・3条、自然法則、ソフトウェア発明。

\paragraph{考え方・正答を選ぶ理由（査読待ち初稿）} 4項とも基本定義に沿い4である。エの除外は抽象的な規則・数学・情報の単なる提示についてであり、暗号や算法を使う装置・ソフトウェア全般の特許を禁止するものではない。請求項全体としての自然法則の利用を判断する。発明に該当するだけで直ちに特許になるわけではなく、新規性・進歩性等も別要件である。

\Needspace{4\baselineskip}\paragraph{選択肢の検討（査読待ち初稿）}\mbox{}\par

\noindent\textbf{1：設問条件に不一致} 正／正／正／誤。エの正誤が違う。各記述の定義・条件を先に判定してから表に対応させる。\par

\noindent\textbf{2：設問条件に不一致} 正／誤／誤／正。イ・ウの正誤が違う。各記述の定義・条件を先に判定してから表に対応させる。\par

\noindent\textbf{3：設問条件に不一致} 正／誤／誤／誤。イ・ウ・エの正誤が違う。各記述の定義・条件を先に判定してから表に対応させる。\par

\noindent\textbf{4：設問条件に一致} 正／正／正／正。全項の判断に一致する。\par

\noindent\textbf{5：設問条件に不一致} 誤／正／正／誤。ア・エの正誤が違う。各記述の定義・条件を先に判定してから表に対応させる。\par

\paragraph{再利用する考え方} 技術的創作と単なる規則・情報の区別について、対象・要件・例外・時点を分け、類似用語や無条件の断定を避ける。

\paragraph{共通理論} \hyperref[theory:compliance:ip]{第\ref*{theory:compliance:ip}節}, \hyperref[theory:compliance:invention]{第\ref*{theory:compliance:invention}節}

\par 参考：\cite{ip_basic_act,jpo_patent_utility,jpo_invention_criteria}。

\par 難易度 3（校準初稿）、優先度未評価。

\subsection{令和4年度 問9}
\label{q:r04-reg-apt-09}

\noindent\textbf{原典：} sources/R04/aptitude.pdf、PDF物理8頁。

\textbf{公式正答：} 5（single）。

\paragraph{問題の要旨} 知的財産のうち創作物に関する5例の正誤の組を選ぶ。

\Needspace{4\baselineskip}\paragraph{確認した条件・記述（独自要約）}\mbox{}\par

\noindent\textbf{ア：} 発明（特許法）。\par

\noindent\textbf{イ：} 考案（実用新案法）。\par

\noindent\textbf{ウ：} 意匠（意匠法）。\par

\noindent\textbf{エ：} 著作物（著作権法）。\par

\noindent\textbf{オ：} 営業秘密（不正競争防止法）。\par

\noindent\textbf{並び：} ア／イ／ウ／エ／オ。「産業財産権」の数ではなく創作物の区分。\par

\paragraph{問われる核心} 知的創作物と産業財産権の範囲の違い

\paragraph{必要知識} 発明・考案・意匠・著作物・営業秘密、創作物と標識、登録権利と保護利益。

\paragraph{考え方・正答を選ぶ理由（査読待ち初稿）} 5例は全て知的創作物の保護に関わるため5である。営業秘密は秘密管理性・有用性・非公知性を満たす情報で、不正競争防止法が保護する。特許庁に登録して生じる4種類の産業財産権に限定すると、著作物や営業秘密を誤って除外してしまう。

\Needspace{4\baselineskip}\paragraph{選択肢の検討（査読待ち初稿）}\mbox{}\par

\noindent\textbf{1：設問条件に不一致} 正／誤／正／正／正。イの正誤が違う。各記述の定義・条件を先に判定してから表に対応させる。\par

\noindent\textbf{2：設問条件に不一致} 正／正／誤／正／正。ウの正誤が違う。各記述の定義・条件を先に判定してから表に対応させる。\par

\noindent\textbf{3：設問条件に不一致} 正／正／正／誤／正。エの正誤が違う。各記述の定義・条件を先に判定してから表に対応させる。\par

\noindent\textbf{4：設問条件に不一致} 正／正／正／正／誤。オの正誤が違う。各記述の定義・条件を先に判定してから表に対応させる。\par

\noindent\textbf{5：設問条件に一致} 正／正／正／正／正。全項の判断に一致する。\par

\paragraph{再利用する考え方} 知的創作物と産業財産権の範囲の違いについて、対象・要件・例外・時点を分け、類似用語や無条件の断定を避ける。

\paragraph{共通理論} \hyperref[theory:compliance:ip]{第\ref*{theory:compliance:ip}節}, \hyperref[theory:compliance:secret]{第\ref*{theory:compliance:secret}節}

\par 参考：\cite{ip_basic_act,jpo_ip_overview,meti_trade_secret}。

\paragraph{関連問題} \hyperref[q:r05-reg-apt-04]{R05-REG-APT-04}

\par 難易度 3（校準初稿）、優先度未評価。

\subsection{令和5年度 問4}
\label{q:r05-reg-apt-04}

\noindent\textbf{原典：} sources/R05/aptitude.pdf、PDF物理4頁。

\textbf{公式正答：} 2（single）。

\paragraph{問題の要旨} 産業財産権に含まれる5種類の権利の正誤の組を選ぶ。

\Needspace{4\baselineskip}\paragraph{確認した条件・記述（独自要約）}\mbox{}\par

\noindent\textbf{並び：} ア：特許権／イ：実用新案権／ウ：意匠権／エ：商標権／オ：著作権。\par

\paragraph{問われる核心} 産業財産権4種類と広い知的財産権の区別

\paragraph{必要知識} 特許庁所管4権利、著作権、権利の対象と制度の範囲。

\paragraph{考え方・正答を選ぶ理由（査読待ち初稿）} 産業財産権は特許・実用新案・意匠・商標の4種類である。著作権も知的財産権であるが、この4種類に含まれない。正／正／正／正／誤の2を選ぶ。産業に使われる著作物だから産業財産権になるという分類ではない。

\Needspace{4\baselineskip}\paragraph{選択肢の検討（査読待ち初稿）}\mbox{}\par

\noindent\textbf{1：設問条件に不一致} 正／正／正／正／正。オの正誤が違う。各記述の定義・条件を先に判定してから表に対応させる。\par

\noindent\textbf{2：設問条件に一致} 正／正／正／正／誤。全項の判断に一致する。\par

\noindent\textbf{3：設問条件に不一致} 正／正／正／誤／正。エ・オの正誤が違う。各記述の定義・条件を先に判定してから表に対応させる。\par

\noindent\textbf{4：設問条件に不一致} 正／正／誤／正／正。ウ・オの正誤が違う。各記述の定義・条件を先に判定してから表に対応させる。\par

\noindent\textbf{5：設問条件に不一致} 正／誤／正／正／正。イ・オの正誤が違う。各記述の定義・条件を先に判定してから表に対応させる。\par

\paragraph{再利用する考え方} 産業財産権4種類と広い知的財産権の区別について、対象・要件・例外・時点を分け、類似用語や無条件の断定を避ける。

\paragraph{共通理論} \hyperref[theory:compliance:ip]{第\ref*{theory:compliance:ip}節}

\par 参考：\cite{jpo_ip_overview}。

\paragraph{関連問題} \hyperref[q:r04-reg-apt-09]{R04-REG-APT-09}

\par 難易度 3（校準初稿）、優先度未評価。

\subsection{令和6年度 問3}
\label{q:r06-reg-apt-03}

\noindent\textbf{原典：} sources/R06/aptitude.pdf、PDF物理3頁。

\textbf{公式正答：} 1（single）。

\paragraph{問題の要旨} 引用・官公資料・AI学習・技術試験の著作権制限について正誤の組を選ぶ。

\Needspace{4\baselineskip}\paragraph{確認した条件・記述（独自要約）}\mbox{}\par

\noindent\textbf{ア：} 学術論文を無承諾で引用すれば全て侵害となる。\par

\noindent\textbf{イ：} 広く利用すべき官公資料は、転載禁止表示がない等の要件の下で説明材料に転載できる。\par

\noindent\textbf{ウ：} AI開発目的の学習データ収集や開発する第三者への提供は権利制限の対象となる。\par

\noindent\textbf{エ：} 複製向けのカメラ・プリンター・和紙等を開発するための試験的複製は権利制限の対象となる。\par

\noindent\textbf{並び：} ア／イ／ウ／エ。2024年。必要な限度と権利者の利益を不当に害しない条件を適用。\par

\paragraph{問われる核心} 著作権保護と利用を調整する要件付きの権利制限

\paragraph{必要知識} 著作権法32条・30条の4、引用、官公資料、非享受目的、必要限度、不当な利益侵害、2019年施行改正。

\paragraph{考え方・正答を選ぶ理由（査読待ち初稿）} アは32条の適法引用を無視しているので誤り。イは32条2項の資料・刊行物等の要件内、ウ・エは30条の4の非享受目的・必要な限度・不当な利益侵害をしない範囲の例で正しい。誤／正／正／正の1である。AI学習に使うと言えば常に自由になるわけではなく、作品の表現を出力させる享受目的が併存する利用や解析用データベースの市場を害する利用は別に検討する。生成した出力の利用も学習段階とは別である。

\Needspace{4\baselineskip}\paragraph{選択肢の検討（査読待ち初稿）}\mbox{}\par

\noindent\textbf{1：設問条件に一致} 誤／正／正／正。全項の判断に一致する。\par

\noindent\textbf{2：設問条件に不一致} 正／誤／正／正。ア・イの正誤が違う。各記述の定義・条件を先に判定してから表に対応させる。\par

\noindent\textbf{3：設問条件に不一致} 正／正／誤／正。ア・ウの正誤が違う。各記述の定義・条件を先に判定してから表に対応させる。\par

\noindent\textbf{4：設問条件に不一致} 正／正／正／誤。ア・エの正誤が違う。各記述の定義・条件を先に判定してから表に対応させる。\par

\noindent\textbf{5：設問条件に不一致} 正／正／正／正。アの正誤が違う。各記述の定義・条件を先に判定してから表に対応させる。\par

\paragraph{再利用する考え方} 著作権保護と利用を調整する要件付きの権利制限について、対象・要件・例外・時点を分け、類似用語や無条件の断定を避ける。

\paragraph{共通理論} \hyperref[theory:compliance:citation]{第\ref*{theory:compliance:citation}節}, \hyperref[theory:compliance:nonenjoyment]{第\ref*{theory:compliance:nonenjoyment}節}

\par 参考：\cite{bunka_citation_faq,bunka_amend2018,bunka_ai2024,bunka_text2026}。

\paragraph{関連問題} \hyperref[q:h27-reg-apt-12]{H27-REG-APT-12}

\par 難易度 3（校準初稿）、優先度未評価。

\section{公正取引・競争法・金融市場}

\subsection{平成27年度 問10}
\label{q:h27-reg-apt-10}

\noindent\textbf{原典：} sources/H27/aptitude.pdf、PDF物理10頁。

\textbf{公式正答：} 5（single）。

\paragraph{問題の要旨} 公共・民間入札の契約方式と談合への法的対応から、不適切な説明を選ぶ。

\Needspace{4\baselineskip}\paragraph{確認した条件・記述（独自要約）}\mbox{}\par

\noindent\textbf{時点：} 平成27年の制度。\par

\noindent\textbf{方向：} 最も不適切なもの。\par

\paragraph{問われる核心} 公共契約方式と発注主体に依存しない競争法規制

\paragraph{必要知識} 一般競争・随意契約、独占禁止法、不当な取引制限、行政・刑事責任、官製談合。

\paragraph{考え方・正答を選ぶ理由（査読待ち初稿）} 発注者が公的か民間かという区別と、事業者間で競争を制限したかという区別は別である。民間入札なら適法という5の断定が誤りである。公共契約の例外、競争法上の違反、職員の関与の責任も同じ制度へまとめない。

\Needspace{4\baselineskip}\paragraph{選択肢の検討（査読待ち初稿）}\mbox{}\par

\noindent\textbf{1：設問条件に不一致} 公共調達は一般競争入札が基本である。。国の会計法上、一般競争を原則とする。地方公共団体等の具体的根拠は各制度に従う。\par

\noindent\textbf{2：設問条件に不一致} 法定条件では随意契約が認められる。。性質・目的が競争を許さない、緊急等の法定要件があり、自由に相手を選ぶ例外ではない。\par

\noindent\textbf{3：設問条件に不一致} 入札談合に行政処分と刑事罰が科され得る。。独占禁止法上の措置や課徴金、刑事処罰の対象となり得る。ただし各行為・責任の要件があり、自動的に全参加者へ同じ制裁が及ぶとはしない。\par

\noindent\textbf{4：設問条件に不一致} 官製談合への職員の関与に刑事罰の制度がある。。官製談合防止制度には一定の職員の入札等妨害への刑事罰があり、2015年以前に施行されている。\par

\noindent\textbf{5：設問条件に一致} 民間事業者による入札の談合は独占禁止法等に抵触しない。。独占禁止法は民間調達の競争制限を一律に除外しない。民間の受注調整にも実際の法的措置がある。\par

\paragraph{再利用する考え方} 公共契約方式と発注主体に依存しない競争法規制について、対象・要件・例外・時点を分け、類似用語や無条件の断定を避ける。

\paragraph{共通理論} \hyperref[theory:compliance:procurement]{第\ref*{theory:compliance:procurement}節}, \hyperref[theory:compliance:competition]{第\ref*{theory:compliance:competition}節}

\par 参考：\cite{mof_procurement2006,jftc_rules,jftc_kansei,jftc_private2018}。

\par 難易度 3（校準初稿）、優先度未評価。

\subsection{令和4年度 問12}
\label{q:r04-reg-apt-12}

\noindent\textbf{原典：} sources/R04/aptitude.pdf、PDF物理11頁。

\textbf{公式正答：} 5（single）。

\paragraph{問題の要旨} 公正取引を損なう4行為と名称について正誤の組を選ぶ。

\Needspace{4\baselineskip}\paragraph{確認した条件・記述（独自要約）}\mbox{}\par

\noindent\textbf{ア：} ア：公共入札で事前に受注者等を決める行為をインサイダー取引という。\par

\noindent\textbf{イ：} イ：自然な需給に見せて人為的に価格を動かし利益を図る行為を相場操縦という。\par

\noindent\textbf{ウ：} ウ：事業者間で商品価格や生産・販売数量を共同決定する行為を談合という。\par

\noindent\textbf{エ：} エ：会社関係者等が未公表の重要事実に関係する株式等の取引をする行為をカルテルという。\par

\paragraph{問われる核心} 競争制限と金融市場の公正を損なう行為の識別

\paragraph{必要知識} カルテル、入札談合、インサイダー取引、相場操縦、規制する行為と法の違い。

\paragraph{考え方・正答を選ぶ理由（査読待ち初稿）} 入札の受注者・価格の事前調整は入札談合、商品の価格・数量の共同決定はカルテル、未公表の重要事実に基づく関係者等の規制はインサイダー取引、人為的な価格形成で他者を誘引する行為は相場操縦である。誤／正／誤／誤の5。なお、入札談合は競争法上カルテルの典型例でもある。用語の出題上の区別を、法的な概念が全く重ならないという主張にしない。

\Needspace{4\baselineskip}\paragraph{選択肢の検討（査読待ち初稿）}\mbox{}\par

\noindent\textbf{1：設問条件に不一致} 正／誤／正／正。ア・イ・ウ・エの正誤が違う。各記述の定義・条件を先に判定してから表に対応させる。\par

\noindent\textbf{2：設問条件に不一致} 正／正／正／誤。ア・ウの正誤が違う。各記述の定義・条件を先に判定してから表に対応させる。\par

\noindent\textbf{3：設問条件に不一致} 誤／正／誤／正。エの正誤が違う。各記述の定義・条件を先に判定してから表に対応させる。\par

\noindent\textbf{4：設問条件に不一致} 正／誤／誤／正。ア・イ・エの正誤が違う。各記述の定義・条件を先に判定してから表に対応させる。\par

\noindent\textbf{5：設問条件に一致} 誤／正／誤／誤。全項の判断に一致する。\par

\paragraph{再利用する考え方} 競争制限と金融市場の公正を損なう行為の識別について、対象・要件・例外・時点を分け、類似用語や無条件の断定を避ける。

\paragraph{共通理論} \hyperref[theory:compliance:competition]{第\ref*{theory:compliance:competition}節}, \hyperref[theory:compliance:market]{第\ref*{theory:compliance:market}節}

\par 参考：\cite{jftc_rules,jftc_bid_cartel2021,fsa_market2004,fsa_information}。

\paragraph{関連問題} \hyperref[q:r07-reg-apt-11]{R07-REG-APT-11}

\par 難易度 3（校準初稿）、優先度未評価。

\subsection{令和7年度 問11}
\label{q:r07-reg-apt-11}

\noindent\textbf{原典：} sources/R07/aptitude.pdf、PDF物理11頁。

\textbf{公式正答：} 4（single）。

\paragraph{問題の要旨} 公正取引を損なう4行為と名称について適切な記述の数を選ぶ。

\Needspace{4\baselineskip}\paragraph{確認した条件・記述（独自要約）}\mbox{}\par

\noindent\textbf{ア：} ア：事業者間で商品価格や生産・販売数量を共同決定する行為を談合という。\par

\noindent\textbf{イ：} イ：会社関係者等が未公表の重要事実に関係する株式等の取引をする行為をカルテルという。\par

\noindent\textbf{ウ：} ウ：自然な需給に見せて人為的に価格を動かし利益を図る行為を相場操縦という。\par

\noindent\textbf{エ：} エ：公共入札で事前に受注者等を決める行為をインサイダー取引という。\par

\paragraph{問われる核心} 競争制限と金融市場の公正を損なう行為の識別

\paragraph{必要知識} カルテル、入札談合、インサイダー取引、相場操縦、規制する行為と法の違い。

\paragraph{考え方・正答を選ぶ理由（査読待ち初稿）} 入札の受注者・価格の事前調整は入札談合、商品の価格・数量の共同決定はカルテル、未公表の重要事実に基づく関係者等の規制はインサイダー取引、人為的な価格形成で他者を誘引する行為は相場操縦である。ウだけ正しく1項なので4。なお、入札談合は競争法上カルテルの典型例でもある。用語の出題上の区別を、法的な概念が全く重ならないという主張にしない。

\Needspace{4\baselineskip}\paragraph{選択肢の検討（査読待ち初稿）}\mbox{}\par

\noindent\textbf{1：設問条件に不一致} 4。適切なのはウだけで、4項ではない。\par

\noindent\textbf{2：設問条件に不一致} 3。アはカルテル、イはインサイダー、エは入札談合なので3項ではない。\par

\noindent\textbf{3：設問条件に不一致} 2。ウ以外は行為と名称が合わず、2項ではない。\par

\noindent\textbf{4：設問条件に一致} 1。相場操縦のウだけが適切で、1項である。\par

\noindent\textbf{5：設問条件に不一致} 0。ウの説明が適切なので0項ではない。\par

\paragraph{再利用する考え方} 競争制限と金融市場の公正を損なう行為の識別について、対象・要件・例外・時点を分け、類似用語や無条件の断定を避ける。

\paragraph{共通理論} \hyperref[theory:compliance:competition]{第\ref*{theory:compliance:competition}節}, \hyperref[theory:compliance:market]{第\ref*{theory:compliance:market}節}

\par 参考：\cite{jftc_rules,jftc_bid_cartel2021,fsa_market2004,fsa_information}。

\paragraph{関連問題} \hyperref[q:r04-reg-apt-12]{R04-REG-APT-12}

\par 難易度 3（校準初稿）、優先度未評価。

\section{引用・官公資料転載・翻訳の条件}

\subsection{平成25年度 問13}
\label{q:h25-reg-apt-13}

\noindent\textbf{原典：} sources/H25/aptitude.pdf、PDF物理14頁。

\textbf{公式正答：} 1（single）。

\par\textbf{検討未確定：} 原典画像の条件と5候補は確認済み。現在取得条文で引用・翻訳の一般原理を照合したが、2013・2014年の当時版全文の独立照合は未完。現在の47条の6を過去の条番号へ代用しない。

\paragraph{問題の要旨} 出所明示、官公資料、外国語論文、許諾なし引用について過度な一般化を判定する。

\Needspace{4\baselineskip}\paragraph{確認した条件・記述（独自要約）}\mbox{}\par

\noindent\textbf{判断対象：} 広義知的財産・産業財産権・引用・営業秘密を区別する。\par

\noindent\textbf{原典照合：} 15問の原典画像について問題幹、否定、表の列順、全5候補を確認。\par

\noindent\textbf{法の時点：} 出題時の制度を保持し、説明の取得基準日は2026-10-12とする。\par

\paragraph{問われる核心} 出所明示、官公資料、外国語論文、許諾なし引用について過度な一般化を判定する。

\paragraph{必要知識} 正誤表、公表済み著作物、公正な慣行、正当な範囲、出所、官公資料、翻訳。

\paragraph{考え方・正答を選ぶ理由（査読待ち初稿）} アは出所を明示すればどんな利用量・目的でも許されるとし、引用の要件を欠くため誤である。イは官公資料を無条件に転載できるとするが、32条2項等には対象・目的・禁止表示の条件があるため誤である。ウは適法な引用の翻訳利用を排除するため誤、エは引用の権利制限を否定するため誤である。全て誤の1が独立判断であり公式と一致する。当時版全文の照合は残し、現在条文の番号を出題時点へ遡及しない。

\Needspace{4\baselineskip}\paragraph{選択肢の検討（査読待ち初稿）}\mbox{}\par

\noindent\textbf{1：独立判断で適切} 誤・誤・誤・誤。4文とも条件を消した断定であり、全て誤となる。\par

\noindent\textbf{2：不適切} 誤・正・誤・誤。官公資料も対象と目的、禁止表示などの条件があり、イは正ではない。\par

\noindent\textbf{3：不適切} 正・正・誤・誤。出典だけで自由利用できず、官公資料も無条件ではないためア・イが違う。\par

\noindent\textbf{4：不適切} 正・正・正・誤。ア・イに加え、引用の翻訳を常に禁止するウも違う。\par

\noindent\textbf{5：不適切} 正・正・正・正。ア〜ウに加えて、適法な引用まで例外なく違法とするエも違う。\par

\paragraph{再利用する考え方} 保護する対象、制度の範囲、成立要件、許諾なし利用の条件を分ける。典型例の名称から成立要件を省略せず、正誤表の行と対象数を対応させる。

\paragraph{共通理論} \hyperref[theory:ip:quotation]{第\ref*{theory:ip:quotation}節}, \hyperref[theory:ip:government]{第\ref*{theory:ip:government}節}, \hyperref[theory:ip:translation]{第\ref*{theory:ip:translation}節}, \hyperref[theory:ip:historical]{第\ref*{theory:ip:historical}節}

\par 参考：\cite{ip_copyright_statute}。

\paragraph{関連問題} \hyperref[q:h26-reg-apt-14]{H26-REG-APT-14}

\par 難易度 3（校準初稿）、優先度未評価。

\section{著作権の保護対象と無方式主義}

\subsection{平成24年度 問9}
\label{q:h24-reg-apt-09}

\noindent\textbf{原典：} sources/H24/aptitude.pdf、PDF物理8頁。

\textbf{公式正答：} 2（single）。

\paragraph{問題の要旨} 知的財産の趣旨、産業財産権の範囲、著作権の発生と公開利用を対応させる。

\Needspace{4\baselineskip}\paragraph{確認した条件・記述（独自要約）}\mbox{}\par

\noindent\textbf{判断対象：} 広義知的財産・産業財産権・引用・営業秘密を区別する。\par

\noindent\textbf{原典照合：} 15問の原典画像について問題幹、否定、表の列順、全5候補を確認。\par

\noindent\textbf{法の時点：} 出題時の制度を保持し、説明の取得基準日は2026-10-12とする。\par

\paragraph{問われる核心} 知的財産の趣旨、産業財産権の範囲、著作権の発生と公開利用を対応させる。

\paragraph{必要知識} 保護対象、4種類、無方式主義、同一性保持、公衆送信、三群の対応。

\paragraph{考え方・正答を選ぶ理由（査読待ち初稿）} 第1群は創作者の保護を通じ創造活動を支えるロ、第2群は特許・実用新案・意匠・商標のイ、第3群は無許諾等の写真掲載を問題にするハであり2となる。株式の権利と紙の株券という媒体の区別は、電子化で新たに知的財産権になることを意味しない。著作権は創作的な表現を保護し、申請・審査なしに生じる。歌唱する者の持ち歌という事情は、作詞者の同意なしに歌詞を変える一般的権限を与えない。趣味の公開Webは私的複製の範囲ではなく、許諾や適法な引用の根拠を別に確認する。

\Needspace{4\baselineskip}\paragraph{選択肢の検討（査読待ち初稿）}\mbox{}\par

\noindent\textbf{1：不適切} イ・イ・ロ。第1群の株券の電子化を知的財産の成立説明とすることと、第3群の歌詞を自由変更する説明が不適切である。\par

\noindent\textbf{2：独立判断で適切} ロ・イ・ハ。創作への誘因、4種類の範囲、許諾等なく公開Webへ写真集の写真を載せる問題が対応する。\par

\noindent\textbf{3：不適切} ロ・ロ・イ。出版権を産業財産権とする第2群と、特許庁の審査で著作権が発生するとする第3群が違う。\par

\noindent\textbf{4：不適切} ハ・ロ・イ。特許の対象を表現、著作権の対象をアイデアへ逆転し、出版権の分類と著作権の発生も違う。\par

\noindent\textbf{5：不適切} ハ・ハ・ハ。最初の保護対象を逆転し、営業秘密・回路配置利用権を産業財産権に含める点も違う。\par

\paragraph{再利用する考え方} 保護する対象、制度の範囲、成立要件、許諾なし利用の条件を分ける。典型例の名称から成立要件を省略せず、正誤表の行と対象数を対応させる。

\paragraph{共通理論} \hyperref[theory:ip:scope]{第\ref*{theory:ip:scope}節}, \hyperref[theory:ip:copyright]{第\ref*{theory:ip:copyright}節}, \hyperref[theory:ip:personality]{第\ref*{theory:ip:personality}節}, \hyperref[theory:ip:public_web]{第\ref*{theory:ip:public_web}節}

\par 参考：\cite{ip_industrial_four,ip_copyright_statute}。

\par 難易度 2（校準初稿）、優先度未評価。

\subsection{平成26年度 問14}
\label{q:h26-reg-apt-14}

\noindent\textbf{原典：} sources/H26/aptitude.pdf、PDF物理14頁。

\textbf{公式正答：} 4（single）。

\par\textbf{検討未確定：} 原典画像の条件と5候補は確認済み。現在取得条文で引用・翻訳の一般原理を照合したが、2013・2014年の当時版全文の独立照合は未完。現在の47条の6を過去の条番号へ代用しない。

\paragraph{問題の要旨} 創作性、著作者人格権、引用要件、外国語引用を判定する。

\Needspace{4\baselineskip}\paragraph{確認した条件・記述（独自要約）}\mbox{}\par

\noindent\textbf{判断対象：} 広義知的財産・産業財産権・引用・営業秘密を区別する。\par

\noindent\textbf{原典照合：} 15問の原典画像について問題幹、否定、表の列順、全5候補を確認。\par

\noindent\textbf{法の時点：} 出題時の制度を保持し、説明の取得基準日は2026-10-12とする。\par

\paragraph{問われる核心} 創作性、著作者人格権、引用要件、外国語引用を判定する。

\paragraph{必要知識} 正誤表、創作的な表現、人格権、引用条件、翻訳利用。

\paragraph{考え方・正答を選ぶ理由（査読待ち初稿）} アは創作性が不要になったという説明で誤、イは公表・氏名表示・同一性保持の人格権を述べて正、ウは公正な慣行と目的上正当な範囲の引用で正、エは外国語のままに限る断定で誤となる。誤・正・正・誤の4を選ぶ。権利の保護と公正な利用を調整する趣旨は、創作性の削除や翻訳の一律禁止を導かない。

\Needspace{4\baselineskip}\paragraph{選択肢の検討（査読待ち初稿）}\mbox{}\par

\noindent\textbf{1：不適切} 正・正・誤・正。創作性不要を正とし、正当な引用の要件を誤とし、翻訳禁止を正とする点が違う。\par

\noindent\textbf{2：不適切} 正・誤・正・誤。創作性は必要であり、人格権の3分類も適切なので前二文が違う。\par

\noindent\textbf{3：不適切} 誤・正・誤・正。引用要件を誤、翻訳禁止を正とする後二文が逆である。\par

\noindent\textbf{4：独立判断で適切} 誤・正・正・誤。創作性必要、人格権の分類、引用の条件、翻訳を一律禁止しないことからこの組合せとなる。\par

\noindent\textbf{5：不適切} 誤・誤・誤・誤。人格権と引用要件を述べたイ・ウは正であり、全て誤ではない。\par

\paragraph{再利用する考え方} 保護する対象、制度の範囲、成立要件、許諾なし利用の条件を分ける。典型例の名称から成立要件を省略せず、正誤表の行と対象数を対応させる。

\paragraph{共通理論} \hyperref[theory:ip:copyright]{第\ref*{theory:ip:copyright}節}, \hyperref[theory:ip:personality]{第\ref*{theory:ip:personality}節}, \hyperref[theory:ip:quotation]{第\ref*{theory:ip:quotation}節}, \hyperref[theory:ip:translation]{第\ref*{theory:ip:translation}節}

\par 参考：\cite{ip_copyright_statute}。

\paragraph{関連問題} \hyperref[q:h25-reg-apt-13]{H25-REG-APT-13}

\par 難易度 3（校準初稿）、優先度未評価。

\section{非享受目的の利用と権利制限の境界}

\subsection{令和元年度再試験 問9}
\label{q:r01-ret-apt-09}

\noindent\textbf{原典：} sources/R01\_retest/aptitude.pdf、PDF物理9頁。

\textbf{公式正答：} 5（single）。

\paragraph{問題の要旨} 2019年施行の非享受目的の利用に該当する六事例の数を求める。

\Needspace{4\baselineskip}\paragraph{確認した条件・記述（独自要約）}\mbox{}\par

\noindent\textbf{判断対象：} 広義知的財産・産業財産権・引用・営業秘密を区別する。\par

\noindent\textbf{原典照合：} 15問の原典画像について問題幹、否定、表の列順、全5候補を確認。\par

\noindent\textbf{法の時点：} 出題時の制度を保持し、説明の取得基準日は2026-10-12とする。\par

\paragraph{問われる核心} 2019年施行の非享受目的の利用に該当する六事例の数を求める。

\paragraph{必要知識} 2019年施行30条の4、非享受目的、必要限度、不当利益侵害、六例の境界。

\paragraph{考え方・正答を選ぶ理由（査読待ち初稿）} 各例の目的を見る。アはAI学習の解析、イはプログラム解析、ウは複製機器や和紙の性能試験、エは送り仮名等の分析、オは写真の表現を鑑賞するのでなく場所情報の3DCG化、カは全文内部複製から書名・著者等を検索する利用である。文化庁の2019年説明資料13〜21物理頁は六例と境界を具体的に示す。該当数6に対応する5を選ぶ。必要限度と不当利益侵害の除外は全例にかかり、3DCGに元写真の創作的表現を承継して享受させる場合、検索結果に本文を表示する場合、AIの享受目的が併存する場合は個別に区別する。

\Needspace{4\baselineskip}\paragraph{選択肢の検討（査読待ち初稿）}\mbox{}\par

\noindent\textbf{1：不適切} 2事例。AI、解析、性能試験、表記調査、場所再現、検索の全6例が該当例であり2ではない。\par

\noindent\textbf{2：不適切} 3事例。目的を表現の鑑賞と混同して三例を外す根拠はない。\par

\noindent\textbf{3：不適切} 4事例。場所の3DCGと検索も非享受の例になり得るため、四例に限定しない。\par

\noindent\textbf{4：不適切} 5事例。六例中一つを常に除く判断は文化庁の例示に対応しない。\par

\noindent\textbf{5：独立判断で適切} 6事例。原典と同じ六例が2019年文化庁資料にあり、条件を満たす非享受利用の例として6となる。\par

\paragraph{再利用する考え方} 保護する対象、制度の範囲、成立要件、許諾なし利用の条件を分ける。典型例の名称から成立要件を省略せず、正誤表の行と対象数を対応させる。

\paragraph{共通理論} \hyperref[theory:ip:nonenjoyment]{第\ref*{theory:ip:nonenjoyment}節}, \hyperref[theory:ip:six_cases]{第\ref*{theory:ip:six_cases}節}, \hyperref[theory:ip:historical]{第\ref*{theory:ip:historical}節}

\par 参考：\cite{ip_flexible2019,bunka_amend2018,ip_copyright_statute}。

\par 難易度 2（校準初稿）、優先度未評価。

\section{発明の定義と特許の要件}

\subsection{令和元年度再試験 問8}
\label{q:r01-ret-apt-08}

\noindent\textbf{原典：} sources/R01\_retest/aptitude.pdf、PDF物理8頁。

\textbf{公式正答：} 4（single）。

\paragraph{問題の要旨} 発明の定義、公益、産業上利用、新規性と公開の要件を対応させる。

\Needspace{4\baselineskip}\paragraph{確認した条件・記述（独自要約）}\mbox{}\par

\noindent\textbf{判断対象：} 広義知的財産・産業財産権・引用・営業秘密を区別する。\par

\noindent\textbf{原典照合：} 15問の原典画像について問題幹、否定、表の列順、全5候補を確認。\par

\noindent\textbf{法の時点：} 出題時の制度を保持し、説明の取得基準日は2026-10-12とする。\par

\paragraph{問われる核心} 発明の定義、公益、産業上利用、新規性と公開の要件を対応させる。

\paragraph{必要知識} 正誤表、発明該当性、産業上利用、新規性、公序良俗、新規性喪失例外。

\paragraph{考え方・正答を選ぶ理由（査読待ち初稿）} アは特許法2条の発明定義、イは32条の不特許事由に反しないこと、ウは29条の産業上利用可能性で正である。エの既に刊行物等で発表されていることは取得要件ではないため誤で、4となる。出願前公開は通常新規性を損なうが、30条の例外には期限と手続がある。公開なら例外なく不可という説明へ置き換えない。

\Needspace{4\baselineskip}\paragraph{選択肢の検討（査読待ち初稿）}\mbox{}\par

\noindent\textbf{1：不適切} 誤・正・正・誤。アの自然法則を利用した高度な技術的思想の創作は定義どおりで正である。\par

\noindent\textbf{2：不適切} 正・誤・正・正。公益の除外を誤とし、出願前公開を必須にするエを正とする点が違う。\par

\noindent\textbf{3：不適切} 誤・正・誤・正。発明の定義と産業上利用可能性を誤とし、公開を必須にする点も違う。\par

\noindent\textbf{4：独立判断で適切} 正・正・正・誤。ア・イ・ウは要件に対応し、エの既発表を要求する文だけが誤である。\par

\noindent\textbf{5：不適切} 正・正・誤・誤。産業上利用できることは要件なのでウは正である。\par

\paragraph{再利用する考え方} 保護する対象、制度の範囲、成立要件、許諾なし利用の条件を分ける。典型例の名称から成立要件を省略せず、正誤表の行と対象数を対応させる。

\paragraph{共通理論} \hyperref[theory:ip:patent]{第\ref*{theory:ip:patent}節}, \hyperref[theory:ip:novelty]{第\ref*{theory:ip:novelty}節}

\par 参考：\cite{ip_patent_statute,ip_novelty_grace2018}。

\par 難易度 2（校準初稿）、優先度未評価。

\section{知的財産の広い範囲と保護形態}

\subsection{平成29年度 問10}
\label{q:h29-reg-apt-10}

\noindent\textbf{原典：} sources/H29/aptitude.pdf、PDF物理10頁。

\textbf{公式正答：} 1（single）。

\paragraph{問題の要旨} 特許、実用新案、意匠、著作権、営業秘密を広義の知的財産へ分類する。

\Needspace{4\baselineskip}\paragraph{確認した条件・記述（独自要約）}\mbox{}\par

\noindent\textbf{判断対象：} 広義知的財産・産業財産権・引用・営業秘密を区別する。\par

\noindent\textbf{原典照合：} 15問の原典画像について問題幹、否定、表の列順、全5候補を確認。\par

\noindent\textbf{法の時点：} 出題時の制度を保持し、説明の取得基準日は2026-10-12とする。\par

\paragraph{問われる核心} 特許、実用新案、意匠、著作権、営業秘密を広義の知的財産へ分類する。

\paragraph{必要知識} 正誤表、広義の知的財産、保護対象と権利形態の差。

\paragraph{考え方・正答を選ぶ理由（査読待ち初稿）} この問題の範囲は4種類の産業財産権に限定されない。特許・実用新案・意匠に加え、著作権と不正競争防止法の保護対象となる営業秘密も広い知的財産の体系に入る。5文とも正の1となる。営業秘密の保護を、登録された特許と同じ排他的権利と説明しない。

\Needspace{4\baselineskip}\paragraph{選択肢の検討（査読待ち初稿）}\mbox{}\par

\noindent\textbf{1：独立判断で適切} 正・正・正・正・正。5対象はいずれも知的財産の保護体系に入る。\par

\noindent\textbf{2：不適切} 正・正・正・正・誤。営業秘密も法律上保護される利益を含む広い定義に入る。\par

\noindent\textbf{3：不適切} 正・正・正・誤・正。著作権は産業財産権ではないが知的財産である。\par

\noindent\textbf{4：不適切} 正・正・誤・正・正。意匠権はデザインを保護する知的財産である。\par

\noindent\textbf{5：不適切} 正・誤・正・正・正。実用新案権は考案を保護する知的財産である。\par

\paragraph{再利用する考え方} 保護する対象、制度の範囲、成立要件、許諾なし利用の条件を分ける。典型例の名称から成立要件を省略せず、正誤表の行と対象数を対応させる。

\paragraph{共通理論} \hyperref[theory:ip:scope]{第\ref*{theory:ip:scope}節}, \hyperref[theory:ip:secret]{第\ref*{theory:ip:secret}節}

\par 参考：\cite{ip_basic_act,ip_broad_scope}。

\paragraph{関連問題} \hyperref[q:h30-reg-apt-06]{H30-REG-APT-06}

\par 難易度 2（校準初稿）、優先度未評価。

\subsection{平成30年度 問6}
\label{q:h30-reg-apt-06}

\noindent\textbf{原典：} sources/H30/aptitude.pdf、PDF物理6頁。

\textbf{公式正答：} 1（single）。

\paragraph{問題の要旨} 九つの保護対象のうち、広義の知的財産権に含まれない数を求める。

\Needspace{4\baselineskip}\paragraph{確認した条件・記述（独自要約）}\mbox{}\par

\noindent\textbf{判断対象：} 広義知的財産・産業財産権・引用・営業秘密を区別する。\par

\noindent\textbf{原典照合：} 15問の原典画像について問題幹、否定、表の列順、全5候補を確認。\par

\noindent\textbf{法の時点：} 出題時の制度を保持し、説明の取得基準日は2026-10-12とする。\par

\paragraph{問われる核心} 九つの保護対象のうち、広義の知的財産権に含まれない数を求める。

\paragraph{必要知識} 広義の定義、9対象、対象数と選択肢番号の区別。

\paragraph{考え方・正答を選ぶ理由（査読待ち初稿）} 9対象を広い知的財産の体系へ位置付けると、含まれない対象は0個である。選択肢番号は0ではなく1となる。特に営業秘密と商号は、特許庁で登録する4種類とは別の保護を持つ。4種類だけを数えて残りを全て知的財産ではないとしない。

\Needspace{4\baselineskip}\paragraph{選択肢の検討（査読待ち初稿）}\mbox{}\par

\noindent\textbf{1：独立判断で適切} 0個。特許、実用新案、意匠、著作権、回路配置、育成者、営業秘密、商標、商号はいずれも広い体系に入る。\par

\noindent\textbf{2：不適切} 1個。範囲外の対象はなく、一つを除く判断は合わない。\par

\noindent\textbf{3：不適切} 2個。4種類の産業財産権だけを対象にする問題ではない。二つを除く根拠はない。\par

\noindent\textbf{4：不適切} 3個。営業秘密・商号等も法律上の保護利益として含むため、三つを除かない。\par

\noindent\textbf{5：不適切} 4個。9対象のうち四つを除く分類は、原典の広い問いと合わない。\par

\paragraph{再利用する考え方} 保護する対象、制度の範囲、成立要件、許諾なし利用の条件を分ける。典型例の名称から成立要件を省略せず、正誤表の行と対象数を対応させる。

\paragraph{共通理論} \hyperref[theory:ip:scope]{第\ref*{theory:ip:scope}節}, \hyperref[theory:ip:count]{第\ref*{theory:ip:count}節}

\par 参考：\cite{ip_basic_act,ip_broad_scope}。

\paragraph{関連問題} \hyperref[q:h29-reg-apt-10]{H29-REG-APT-10}, \hyperref[q:r02-reg-apt-05]{R02-REG-APT-05}

\par 難易度 2（校準初稿）、優先度未評価。

\section{産業財産権4種類と広義知的財産の境界}

\subsection{平成23年度 問9}
\label{q:h23-reg-apt-09}

\noindent\textbf{原典：} sources/H23/aptitude.pdf、PDF物理10頁。

\textbf{公式正答：} 5（single）。

\paragraph{問題の要旨} 産業財産権4種類に含まれない権利を選ぶ。

\Needspace{4\baselineskip}\paragraph{確認した条件・記述（独自要約）}\mbox{}\par

\noindent\textbf{判断対象：} 広義知的財産・産業財産権・引用・営業秘密を区別する。\par

\noindent\textbf{原典照合：} 15問の原典画像について問題幹、否定、表の列順、全5候補を確認。\par

\noindent\textbf{法の時点：} 出題時の制度を保持し、説明の取得基準日は2026-10-12とする。\par

\paragraph{問われる核心} 産業財産権4種類に含まれない権利を選ぶ。

\paragraph{必要知識} 産業財産権4種類、回路配置利用権と広義の知的財産。

\paragraph{考え方・正答を選ぶ理由（査読待ち初稿）} 特許・実用新案・意匠・商標を先に集合として書く。原典5候補のうち、その外にある回路配置利用権が5である。半導体の技術に関係することと、産業財産権という4種類の総称に入ることを分ける。

\Needspace{4\baselineskip}\paragraph{選択肢の検討（査読待ち初稿）}\mbox{}\par

\noindent\textbf{1：対象範囲を検討} 特許権。発明を保護する産業財産権の一つである。\par

\noindent\textbf{2：対象範囲を検討} 実用新案権。物品の形状等の考案を保護する産業財産権である。\par

\noindent\textbf{3：対象範囲を検討} 意匠権。デザインの保護であり4種類に含まれる。\par

\noindent\textbf{4：対象範囲を検討} 商標権。商品・役務の出所標識を保護する4種類の一つである。\par

\noindent\textbf{5：該当しない対象} 回路配置利用権。回路配置の専用制度は知的財産に関係するが産業財産権の4種類には含まれない。\par

\paragraph{再利用する考え方} 保護する対象、制度の範囲、成立要件、許諾なし利用の条件を分ける。典型例の名称から成立要件を省略せず、正誤表の行と対象数を対応させる。

\paragraph{共通理論} \hyperref[theory:ip:scope]{第\ref*{theory:ip:scope}節}, \hyperref[theory:ip:count]{第\ref*{theory:ip:count}節}

\par 参考：\cite{ip_industrial_four,ip_broad_scope}。

\paragraph{関連問題} \hyperref[q:r01-reg-apt-05]{R01-REG-APT-05}, \hyperref[q:r03-reg-apt-13]{R03-REG-APT-13}

\par 難易度 2（校準初稿）、優先度未評価。

\subsection{令和元年度 問5}
\label{q:r01-reg-apt-05}

\noindent\textbf{原典：} sources/R01/aptitude.pdf、PDF物理6頁。

\textbf{公式正答：} 5（single）。

\paragraph{問題の要旨} 産業財産権に含まれない植物新品種の権利を選ぶ。

\Needspace{4\baselineskip}\paragraph{確認した条件・記述（独自要約）}\mbox{}\par

\noindent\textbf{判断対象：} 広義知的財産・産業財産権・引用・営業秘密を区別する。\par

\noindent\textbf{原典照合：} 15問の原典画像について問題幹、否定、表の列順、全5候補を確認。\par

\noindent\textbf{法の時点：} 出題時の制度を保持し、説明の取得基準日は2026-10-12とする。\par

\paragraph{問われる核心} 産業財産権に含まれない植物新品種の権利を選ぶ。

\paragraph{必要知識} 産業財産権4種類、育成者権。

\paragraph{考え方・正答を選ぶ理由（査読待ち初稿）} 4種類の集合と5候補を比較する。特許・実用新案・意匠・商標は含まれ、育成者権が外れるので5である。平成23年と同じ番号でも、今回は回路配置ではなく植物新品種が対象である。

\Needspace{4\baselineskip}\paragraph{選択肢の検討（査読待ち初稿）}\mbox{}\par

\noindent\textbf{1：対象範囲を検討} 特許権。産業財産権の4種類に含まれる。\par

\noindent\textbf{2：対象範囲を検討} 実用新案権。物品の形状等の考案を保護する4種類の一つである。\par

\noindent\textbf{3：対象範囲を検討} 意匠権。デザインを保護する4種類の一つである。\par

\noindent\textbf{4：対象範囲を検討} 商標権。出所を示す標識の4種類の一つである。\par

\noindent\textbf{5：該当しない対象} 育成者権。植物新品種の育成者権は知的財産に関係するが、この4種類には入らない。\par

\paragraph{再利用する考え方} 保護する対象、制度の範囲、成立要件、許諾なし利用の条件を分ける。典型例の名称から成立要件を省略せず、正誤表の行と対象数を対応させる。

\paragraph{共通理論} \hyperref[theory:ip:scope]{第\ref*{theory:ip:scope}節}, \hyperref[theory:ip:count]{第\ref*{theory:ip:count}節}

\par 参考：\cite{ip_industrial_four,ip_broad_scope}。

\paragraph{関連問題} \hyperref[q:h23-reg-apt-09]{H23-REG-APT-09}, \hyperref[q:r07-reg-apt-08]{R07-REG-APT-08}

\par 難易度 2（校準初稿）、優先度未評価。

\subsection{令和2年度 問5}
\label{q:r02-reg-apt-05}

\noindent\textbf{原典：} sources/R02/aptitude.pdf、PDF物理5頁。

\textbf{公式正答：} 3（single）。

\paragraph{問題の要旨} 十対象のうち産業財産権に含まれない数を求める。

\Needspace{4\baselineskip}\paragraph{確認した条件・記述（独自要約）}\mbox{}\par

\noindent\textbf{判断対象：} 広義知的財産・産業財産権・引用・営業秘密を区別する。\par

\noindent\textbf{原典照合：} 15問の原典画像について問題幹、否定、表の列順、全5候補を確認。\par

\noindent\textbf{法の時点：} 出題時の制度を保持し、説明の取得基準日は2026-10-12とする。\par

\paragraph{問われる核心} 十対象のうち産業財産権に含まれない数を求める。

\paragraph{必要知識} 十対象、産業財産権4種類、補集合の個数。

\paragraph{考え方・正答を選ぶ理由（査読待ち初稿）} 特許・実用新案・意匠・商標だけを含まれる側へ置く。他の著作権、回路配置利用権、育成者権、営業秘密、商号、商品等表示はこの4種類の外である。10-4=\allowbreak{}6個に対応する3を選ぶ。平成30年問6の広い知的財産から外れる数と問いの範囲が異なる。

\Needspace{4\baselineskip}\paragraph{選択肢の検討（査読待ち初稿）}\mbox{}\par

\noindent\textbf{1：不適切} 4個。十対象中四つが含まれるので、4は含まれる側の数である。\par

\noindent\textbf{2：不適切} 5個。著作権、回路配置、育成者、営業秘密、商号、商品等表示の六つが外れる。五つではない。\par

\noindent\textbf{3：独立判断で適切} 6個。10から特許・実用新案・意匠・商標の4を引き、6となる。\par

\noindent\textbf{4：不適切} 7個。4種類の一つを誤って除外しなければ7にはならない。\par

\noindent\textbf{5：不適切} 8個。特許・実用新案以外にも意匠・商標が含まれるため、8ではない。\par

\paragraph{再利用する考え方} 保護する対象、制度の範囲、成立要件、許諾なし利用の条件を分ける。典型例の名称から成立要件を省略せず、正誤表の行と対象数を対応させる。

\paragraph{共通理論} \hyperref[theory:ip:scope]{第\ref*{theory:ip:scope}節}, \hyperref[theory:ip:count]{第\ref*{theory:ip:count}節}

\par 参考：\cite{ip_industrial_four,ip_broad_scope}。

\paragraph{関連問題} \hyperref[q:h30-reg-apt-06]{H30-REG-APT-06}

\par 難易度 2（校準初稿）、優先度未評価。

\subsection{令和3年度 問13}
\label{q:r03-reg-apt-13}

\noindent\textbf{原典：} sources/R03/aptitude.pdf、PDF物理13頁。

\textbf{公式正答：} 3（single）。

\paragraph{問題の要旨} 回路配置利用権を産業財産権4種類から区別する。

\Needspace{4\baselineskip}\paragraph{確認した条件・記述（独自要約）}\mbox{}\par

\noindent\textbf{判断対象：} 広義知的財産・産業財産権・引用・営業秘密を区別する。\par

\noindent\textbf{原典照合：} 15問の原典画像について問題幹、否定、表の列順、全5候補を確認。\par

\noindent\textbf{法の時点：} 出題時の制度を保持し、説明の取得基準日は2026-10-12とする。\par

\paragraph{問われる核心} 回路配置利用権を産業財産権4種類から区別する。

\paragraph{必要知識} 産業財産権4種類、候補位置の変更。

\paragraph{考え方・正答を選ぶ理由（査読待ち初稿）} 特許・実用新案・意匠・商標と候補を比較する。回路配置利用権だけが含まれず3を選ぶ。平成23年は同じ対象が5に置かれていた。選択肢の番号でなく、権利の名称と集合への所属を覚える。

\Needspace{4\baselineskip}\paragraph{選択肢の検討（査読待ち初稿）}\mbox{}\par

\noindent\textbf{1：対象範囲を検討} 特許権。発明の権利で産業財産権に含まれる。\par

\noindent\textbf{2：対象範囲を検討} 実用新案権。物品の考案を保護する4種類の一つである。\par

\noindent\textbf{3：該当しない対象} 回路配置利用権。回路配置を保護する専用制度で、産業財産権4種類の外となる。\par

\noindent\textbf{4：対象範囲を検討} 意匠権。デザインを保護する4種類の一つである。\par

\noindent\textbf{5：対象範囲を検討} 商標権。出所標識を保護する4種類の一つである。\par

\paragraph{再利用する考え方} 保護する対象、制度の範囲、成立要件、許諾なし利用の条件を分ける。典型例の名称から成立要件を省略せず、正誤表の行と対象数を対応させる。

\paragraph{共通理論} \hyperref[theory:ip:scope]{第\ref*{theory:ip:scope}節}, \hyperref[theory:ip:count]{第\ref*{theory:ip:count}節}

\par 参考：\cite{ip_industrial_four,ip_broad_scope}。

\paragraph{関連問題} \hyperref[q:h23-reg-apt-09]{H23-REG-APT-09}

\par 難易度 2（校準初稿）、優先度未評価。

\subsection{令和7年度 問8}
\label{q:r07-reg-apt-08}

\noindent\textbf{原典：} sources/R07/aptitude.pdf、PDF物理8頁。

\textbf{公式正答：} 4（single）。

\paragraph{問題の要旨} 七つの権利から産業財産権でない三つの組合せを選ぶ。

\Needspace{4\baselineskip}\paragraph{確認した条件・記述（独自要約）}\mbox{}\par

\noindent\textbf{判断対象：} 広義知的財産・産業財産権・引用・営業秘密を区別する。\par

\noindent\textbf{原典照合：} 15問の原典画像について問題幹、否定、表の列順、全5候補を確認。\par

\noindent\textbf{法の時点：} 出題時の制度を保持し、説明の取得基準日は2026-10-12とする。\par

\paragraph{問われる核心} 七つの権利から産業財産権でない三つの組合せを選ぶ。

\paragraph{必要知識} 七対象の記号対応、産業財産権の補集合、三対象の組合せ。

\paragraph{考え方・正答を選ぶ理由（査読待ち初稿）} 原典の記号対応を保持する。含まれる側はア特許・ウ意匠・エ商標・カ実用新案であり、外れる側はイ著作権・オ回路配置・キ育成者である。後者の組合せ4を選ぶ。過去の単一対象を選ぶ形式から三対象の組合せへ変わっているが、判断の集合は同じである。

\Needspace{4\baselineskip}\paragraph{選択肢の検討（査読待ち初稿）}\mbox{}\par

\noindent\textbf{1：不適切} ア・オ・カ。アの特許とカの実用新案は含まれるので、除外対象の組合せではない。\par

\noindent\textbf{2：不適切} ア・ウ・キ。アの特許とウの意匠が含まれるため不適切である。\par

\noindent\textbf{3：不適切} イ・ウ・エ。ウの意匠とエの商標が含まれるため不適切である。\par

\noindent\textbf{4：独立判断で適切} イ・オ・キ。イの著作権、オの回路配置利用権、キの育成者権は全て4種類の外である。\par

\noindent\textbf{5：不適切} オ・カ・キ。カの実用新案を除外する点が違う。\par

\paragraph{再利用する考え方} 保護する対象、制度の範囲、成立要件、許諾なし利用の条件を分ける。典型例の名称から成立要件を省略せず、正誤表の行と対象数を対応させる。

\paragraph{共通理論} \hyperref[theory:ip:scope]{第\ref*{theory:ip:scope}節}, \hyperref[theory:ip:count]{第\ref*{theory:ip:count}節}

\par 参考：\cite{ip_industrial_four,ip_broad_scope}。

\paragraph{関連問題} \hyperref[q:r01-reg-apt-05]{R01-REG-APT-05}

\par 難易度 2（校準初稿）、優先度未評価。

\section{営業秘密の三要件と有用性}

\subsection{平成30年度 問7}
\label{q:h30-reg-apt-07}

\noindent\textbf{原典：} sources/H30/aptitude.pdf、PDF物理7頁。

\textbf{公式正答：} 5（single）。

\par\textbf{検討未確定：} 全5組合せを画像確認した。未利用・公開・反社会的情報は当時行政資料の索引箇所と条文で照合したが、サイバー攻撃による漏洩の2018年当時の急増を示す定量的資料は未取得。

\paragraph{問題の要旨} 営業秘密の未利用情報、公開情報、漏洩経路、反社会的情報を判定する。

\Needspace{4\baselineskip}\paragraph{確認した条件・記述（独自要約）}\mbox{}\par

\noindent\textbf{判断対象：} 広義知的財産・産業財産権・引用・営業秘密を区別する。\par

\noindent\textbf{原典照合：} 15問の原典画像について問題幹、否定、表の列順、全5候補を確認。\par

\noindent\textbf{法の時点：} 出題時の制度を保持し、説明の取得基準日は2026-10-12とする。\par

\paragraph{問われる核心} 営業秘密の未利用情報、公開情報、漏洩経路、反社会的情報を判定する。

\paragraph{必要知識} 正誤表、三要件、有用性、非公知性、反社会的情報、当時の漏洩動向。

\paragraph{考え方・正答を選ぶ理由（査読待ち初稿）} アは実際に使われていない情報を一律除外するため誤、イは公開内容の非公知性を否定するので正、エは違法行為の情報を営業秘密とするため誤である。ウの列挙された経路は実際に問題となり得る。原典の組合せは誤・正・正・誤の5で公式と一致するが、ウの2018年当時の「急増」という量的主張を確認した証拠は取得していない。経路の存在と時間的増加を分け、後者の未確認を保持する。

\Needspace{4\baselineskip}\paragraph{選択肢の検討（査読待ち初稿）}\mbox{}\par

\noindent\textbf{1：不適切} 正・正・正・誤。未利用でも有用性があり得るのでアを正としない。\par

\noindent\textbf{2：不適切} 誤・正・誤・誤。ウの経路は存在するが当時の急増系列は未確認であり、原典の真扱いと独立の量的確認を分ける。\par

\noindent\textbf{3：不適切} 正・正・誤・正。未利用情報を一律排除し、反社会的情報を保護する判断が違う。\par

\noindent\textbf{4：不適切} 誤・誤・正・正。公知情報を秘密に含め、反社会的情報も含める判断が違う。\par

\noindent\textbf{5：独立判断で適切} 誤・正・正・誤。未利用情報の可能性、公知の除外、複数漏洩経路、反社会的情報の除外に対応する原典の組合せである。\par

\paragraph{再利用する考え方} 保護する対象、制度の範囲、成立要件、許諾なし利用の条件を分ける。典型例の名称から成立要件を省略せず、正誤表の行と対象数を対応させる。

\paragraph{共通理論} \hyperref[theory:ip:secret]{第\ref*{theory:ip:secret}節}, \hyperref[theory:ip:usefulness]{第\ref*{theory:ip:usefulness}節}, \hyperref[theory:ip:secrecy]{第\ref*{theory:ip:secrecy}節}

\par 参考：\cite{ip_unfair_statute,ip_secret_handbook2017,ip_secret_commentary2018}。

\paragraph{関連問題} \hyperref[q:r03-reg-apt-07]{R03-REG-APT-07}

\par 難易度 3（校準初稿）、優先度未評価。

\subsection{令和2年度 問4}
\label{q:r02-reg-apt-04}

\noindent\textbf{原典：} sources/R02/aptitude.pdf、PDF物理4頁。

\textbf{公式正答：} 2（single）。

\par\textbf{検討未確定：} アは顧客名簿等を典型例として挙げる読みでは真であるが、社内の作成過程だけを十分条件とする一般命題では三要件が欠ける。原典の省略条件を補った確定本文にしない。

\paragraph{問題の要旨} 営業秘密の典型例、反社会的情報、公知情報、三要件の論理を判定する。

\Needspace{4\baselineskip}\paragraph{確認した条件・記述（独自要約）}\mbox{}\par

\noindent\textbf{判断対象：} 広義知的財産・産業財産権・引用・営業秘密を区別する。\par

\noindent\textbf{原典照合：} 15問の原典画像について問題幹、否定、表の列順、全5候補を確認。\par

\noindent\textbf{法の時点：} 出題時の制度を保持し、説明の取得基準日は2026-10-12とする。\par

\paragraph{問われる核心} 営業秘密の典型例、反社会的情報、公知情報、三要件の論理を判定する。

\paragraph{必要知識} 正誤表、三要件の論理積、典型例と十分条件、反社会的情報、公知性。

\paragraph{考え方・正答を選ぶ理由（査読待ち初稿）} アは資料名だけでなく秘密管理性・有用性・非公知性が必要である。これらを満たす顧客名簿等の典型例として読むという条件を付け、独立初回は正とした。イは有害物質の違法放出を守る秘密へ含めるため誤、ウは公開情報を秘密から外すので正、エは三要件のいずれか一つでよいとするため誤である。条件付きの正・誤・正・誤は2で公式と一致する。アを無条件の一般命題として読めば偽となり、誤・誤・正・誤の行は原典候補にない。この省略を解消したことにはしない。

\Needspace{4\baselineskip}\paragraph{選択肢の検討（査読待ち初稿）}\mbox{}\par

\noindent\textbf{1：不適切} 正・正・誤・誤。違法な河川放出の情報まで保護するイを正とし、公知の除外を誤とする点が違う。\par

\noindent\textbf{2：独立判断で適切} 正・誤・正・誤。アを三要件を満たす典型例と読む場合の組合せであり、公式2と対応する。アの省略は残る。\par

\noindent\textbf{3：不適切} 誤・誤・正・正。どれか一要件で足りるエを正とするのが違う。アを厳密に偽と読む場合も誤・誤・正・誤で、この候補ではない。\par

\noindent\textbf{4：不適切} 誤・正・誤・正。反社会的情報の保護、公知の扱い、一要件十分という三点が違う。\par

\noindent\textbf{5：不適切} 正・誤・正・正。アの条件付き典型例は読めるが、エは三要件の同時成立なので正にできない。\par

\paragraph{再利用する考え方} 保護する対象、制度の範囲、成立要件、許諾なし利用の条件を分ける。典型例の名称から成立要件を省略せず、正誤表の行と対象数を対応させる。

\paragraph{共通理論} \hyperref[theory:ip:secret]{第\ref*{theory:ip:secret}節}, \hyperref[theory:ip:usefulness]{第\ref*{theory:ip:usefulness}節}, \hyperref[theory:ip:secret_ambiguity]{第\ref*{theory:ip:secret_ambiguity}節}

\par 参考：\cite{ip_unfair_statute,ip_secret_handbook2017}。

\paragraph{関連問題} \hyperref[q:r03-reg-apt-07]{R03-REG-APT-07}

\par 難易度 3（校準初稿）、優先度未評価。

\subsection{令和3年度 問7}
\label{q:r03-reg-apt-07}

\noindent\textbf{原典：} sources/R03/aptitude.pdf、PDF物理7頁。

\textbf{公式正答：} 5（single）。

\par\textbf{独立解答と公式正答に相違あり：} 独立解答 4。アは顧客名簿等を典型例として挙げる読みでは真であるが、社内の作成過程だけを十分条件とする一般命題では三要件が欠ける。原典の省略条件を補った確定本文にしない。独立初回4と公式5の差異が残る。

\paragraph{問題の要旨} 営業秘密の三要件と、顧客名簿を典型例とする省略表現を検討する。

\Needspace{4\baselineskip}\paragraph{確認した条件・記述（独自要約）}\mbox{}\par

\noindent\textbf{判断対象：} 広義知的財産・産業財産権・引用・営業秘密を区別する。\par

\noindent\textbf{原典照合：} 15問の原典画像について問題幹、否定、表の列順、全5候補を確認。\par

\noindent\textbf{法の時点：} 出題時の制度を保持し、説明の取得基準日は2026-10-12とする。\par

\paragraph{問われる核心} 営業秘密の三要件と、顧客名簿を典型例とする省略表現を検討する。

\paragraph{必要知識} 正誤表、営業秘密の三要件、一般命題と典型例、独立判断と公式の差。

\paragraph{考え方・正答を選ぶ理由（査読待ち初稿）} イは反社会的情報が有用性を満たさないので正、ウは公開内容が非公知でないので正、エは秘密管理性が三要件の一つなので正である。アは顧客名簿等が社内活動から生まれたことだけで営業秘密になるという一般命題として読むと誤であり、独立初回は4を保存した。公式正答は5であり、三要件を満たす典型例を述べる読みでは全て正になる。公式値は5として保持し、独立初回4を変更せず、二つの読みの差を未解決として追跡する。令和2年問4の同じアにも条件省略がある。

\Needspace{4\baselineskip}\paragraph{選択肢の検討（査読待ち初稿）}\mbox{}\par

\noindent\textbf{1：不適切} 正・正・正・誤。エの秘密管理性は必要要件なので誤ではない。\par

\noindent\textbf{2：不適切} 正・正・誤・正。ウの公開内容は非公知でないので、除外する説明は正である。\par

\noindent\textbf{3：不適切} 正・誤・正・正。イの反社会的情報の除外は行政資料の有用性の説明と整合し正である。\par

\noindent\textbf{4：一般命題としての独立初回判断} 誤・正・正・正。アを無条件の一般命題として読む独立初回の組合せである。資料名と由来だけでは三要件を満たさない。\par

\noindent\textbf{5：典型例としての公式判断} 正・正・正・正。アを三要件を満たす典型例の説明として読む公式の組合せである。原典で省略された条件を明示して区別する。\par

\paragraph{再利用する考え方} 保護する対象、制度の範囲、成立要件、許諾なし利用の条件を分ける。典型例の名称から成立要件を省略せず、正誤表の行と対象数を対応させる。

\paragraph{共通理論} \hyperref[theory:ip:secret]{第\ref*{theory:ip:secret}節}, \hyperref[theory:ip:usefulness]{第\ref*{theory:ip:usefulness}節}, \hyperref[theory:ip:secret_ambiguity]{第\ref*{theory:ip:secret_ambiguity}節}

\par 参考：\cite{ip_unfair_statute,ip_secret_handbook2017}。

\paragraph{関連問題} \hyperref[q:h30-reg-apt-07]{H30-REG-APT-07}, \hyperref[q:r02-reg-apt-04]{R02-REG-APT-04}

\par 難易度 3（校準初稿）、優先度未評価。

\section{社会的責任・安全の標準}

\subsection{令和6年度 問12}
\label{q:r06-reg-apt-12}

\noindent\textbf{原典：} sources/R06/aptitude.pdf、PDF物理12頁。

\textbf{公式正答：} 3（single）。

\paragraph{問題の要旨} ISO・IEC・JISの名称、認証、審査、図の新旧の6説明について適切なものの数を選ぶ。

\Needspace{4\baselineskip}\paragraph{確認した条件・記述（独自要約）}\mbox{}\par

\noindent\textbf{ア：} ISOは電気・通信・電子技術分野を除く産業分野の国際規格を作る。\par

\noindent\textbf{イ：} IECは電気・電子技術の国際規格を作る。\par

\noindent\textbf{ウ：} 2019年に産業標準化法→工業標準化法、日本産業規格→日本工業規格と改称した。\par

\noindent\textbf{エ：} 国登録の民間第三者機関の適合性審査と認証でJISマークを表示できる。\par

\noindent\textbf{オ：} 登録認証機関は少なくとも3年ごと・必要時臨時に審査し、国も必要に応じ立入検査する。\par

\noindent\textbf{カ：} 図の左を旧、右を現行JISマークと説明している。\par

\paragraph{問われる核心} 国際規格の分担・JISの改称・認証と図の新旧

\paragraph{必要知識} ISO、IEC、産業標準化法、2019年改称、登録認証機関、定期審査、2005年新マーク。

\paragraph{考え方・正答を選ぶ理由（査読待ち初稿）} JISCの分担説明からア・イは正しい。2019年7月1日の改称は工業→産業なのでウは誤り。エ・オは認証・審査制度に沿う。図では左の横並びJISが2005年開始の現行、右が旧で、カの説明は逆である。4項が正しく3を選ぶ。初回独立解答2からの修正は公式番号だけで行わず、JISCの現行マーク画像を別に表示して照合し、その画像・ハッシュ・変更理由と初回記録を保存している。2019年の対象拡大と2005年のマーク制度変更を混同しない。

\Needspace{4\baselineskip}\paragraph{選択肢の検討（査読待ち初稿）}\mbox{}\par

\noindent\textbf{1：設問条件に不一致} 6。ウは改称方向が逆、カは図の新旧が逆なので6項ではない。\par

\noindent\textbf{2：設問条件に不一致} 5。初回の独立判断ではカを正として5項としたが、JISC公式現行図と実画像を照合しカも誤りと確定した。\par

\noindent\textbf{3：設問条件に一致} 4。ア・イ・エ・オの4項が適切である。\par

\noindent\textbf{4：設問条件に不一致} 3。名称と図以外の4項は適切なので3項ではない。\par

\noindent\textbf{5：設問条件に不一致} 2。ウとカが誤りであり、適切なのは2項ではなく4項である。\par

\paragraph{再利用する考え方} 国際規格の分担・JISの改称・認証と図の新旧について、対象・要件・例外・時点を分け、類似用語や無条件の断定を避ける。

\paragraph{共通理論} \hyperref[theory:compliance:standards]{第\ref*{theory:compliance:standards}節}

\par 参考：\cite{jisc_isoiec,jisc_certification,jisc_marks,meti_jis2019,mhlw_jis_history}。

\par 難易度 3（校準初稿）、優先度未評価。

\section{公益通報者保護と通報経路}

\subsection{平成26年度 問6}
\label{q:h26-reg-apt-06}

\noindent\textbf{原典：} sources/H26/aptitude.pdf、PDF物理6頁。

\textbf{公式正答：} 5（single）。

\paragraph{問題の要旨} 公共の福祉、公衆への説明、公務員の秘密保持、公益通報の対象時点を判定する。

\Needspace{4\baselineskip}\paragraph{確認した条件・記述（独自要約）}\mbox{}\par

\noindent\textbf{方向：} 最も不適切。\par

\noindent\textbf{時点：} 平成26年当時の公益通報者保護法。\par

\noindent\textbf{範囲：} 通報対象者・目的・対象事実・通報先には別の条件もある。\par

\paragraph{問われる核心} 公益通報の予防的機能と公衆への責任

\paragraph{必要知識} 公益通報の予防的機能と公衆への責任。1〜4の対象と限定を確認する。5は、事故や犯罪が実際に起こるまで待つ必要があるとの誤った限定である。平成26年当時の法の定義には、まさに生じようとしている場合が含まれる。現行法で対象者や通報対象事実が拡大したこととは別に、当時からあるこの点を根拠に5を選ぶ。

\paragraph{考え方・正答を選ぶ理由（査読待ち初稿）} 1〜4の対象と限定を確認する。5は、事故や犯罪が実際に起こるまで待つ必要があるとの誤った限定である。平成26年当時の法の定義には、まさに生じようとしている場合が含まれる。現行法で対象者や通報対象事実が拡大したこととは別に、当時からあるこの点を根拠に5を選ぶ。

\Needspace{4\baselineskip}\paragraph{選択肢の検討（査読待ち初稿）}\mbox{}\par

\noindent\textbf{1：適切} 自由・権利の保持と公共の福祉に適合する私権。。憲法12条と民法1条1項の趣旨を合わせた記述である。権利の無制限な濫用を認めるものではない。\par

\noindent\textbf{2：適切（技術倫理の文脈）} 技術の結果へ十分な情報に基づく同意を与える立場にない公衆。。当該技術に対する情報・選択機会・力の非対称性を扱う。全ての市民の一般的能力を低く評価する意味ではない。\par

\noindent\textbf{3：適切（説明の責任）} 知る権利と技術者による情報開示・説明責任。。判断に必要な情報を説明する責任を述べる。説明すれば必ず同意する、異論を封じてよいとの意味へ拡張しない。\par

\noindent\textbf{4：適切} 公務員の職務上の秘密保持は退職後も続く。。国家公務員法100条は退職後にも適用する。法定の許可等を要する開示の例外があるので、全ての情報を永久に秘匿する意味ではない。\par

\noindent\textbf{5：不適切} 公益通報の保護は犯罪行為が既に起きた場合のみで、起きようとしている場合を含まない。。出題当時も、通報対象事実が生じ、又はまさに生じようとしている旨を含む。未来の事実を一律に除く後半が誤り。\par

\paragraph{再利用する考え方} 対象、条件、時点を先に確認する。公益通報の予防的機能と公衆への責任を具体的な事例へ適用し、技術的原因・法的要件・倫理的判断のどれを問うか区別する。

\paragraph{共通理論} \hyperref[theory:compliance:whistleblower]{第\ref*{theory:compliance:whistleblower}節}, \hyperref[theory:compliance:accountability]{第\ref*{theory:compliance:accountability}節}, \hyperref[theory:ethics:public]{第\ref*{theory:ethics:public}節}

\par 参考：\cite{caa_handbook_pre2022,caa_report_current,jp_constitution,jp_civil,jp_national_public}。

\par 難易度 2（校準初稿）、優先度未評価。

\subsection{平成30年度 問8}
\label{q:h30-reg-apt-08}

\noindent\textbf{原典：} sources/H30/aptitude.pdf、PDF物理8頁。

\textbf{公式正答：} 3（single）。

\paragraph{問題の要旨} 出題当時の公益通報者保護法について、不適切な説明を選ぶ。

\Needspace{4\baselineskip}\paragraph{確認した条件・記述（独自要約）}\mbox{}\par

\noindent\textbf{時点：} 2018年。2020年改正・2022年施行前。\par

\noindent\textbf{方向：} 最も不適切なもの。\par

\paragraph{問われる核心} 公益通報の対象者・通報先別の保護要件と施行時点

\paragraph{必要知識} 公益通報者保護法2・3・7条、不正目的、相当の理由、公務員、内部・行政・外部、追加条件。

\paragraph{考え方・正答を選ぶ理由（査読待ち初稿）} 公務員を一律に除外する3が誤りである。通報先によって保護要件が異なり、その他外部への通報には当時の3条3号の条件を確認する。原典5の「など」は例示なので、内部通報後20日経過等と急迫した生命・身体の危険を全て同時に要求するとは解さない。現在は2022年施行改正で保護対象・体制・要件が更新されている。2025年改正の主要部分は2026年12月1日施行予定であり、2026年10月10日現在の施行法と混ぜない。

\Needspace{4\baselineskip}\paragraph{選択肢の検討（査読待ち初稿）}\mbox{}\par

\noindent\textbf{1：設問条件に不一致} 不正目的でなく通報対象事実を所定の通報先へ知らせる。。不正の目的を排除し、法律が定める対象事実と通報先を満たす必要がある。\par

\noindent\textbf{2：設問条件に不一致} 保護要件を満たす通報者の解雇・不利益取扱いを防ぐ。。通報を理由とする不利益から保護する制度で、あらゆる行為を無条件に免責するものではない。\par

\noindent\textbf{3：設問条件に一致} 保護対象の労働者に公務員は含まれない。。公務員も原則として労働者に含まれるので誤り。公務員への取扱いには法の特則がある。\par

\noindent\textbf{4：設問条件に不一致} 内部・行政・その他外部で保護要件が異なる。。当時、行政への通報では相当の理由、その他外部ではさらに所定の追加条件などが必要であった。\par

\noindent\textbf{5：設問条件に不一致} その他外部への通報には相当の理由と被害防止に必要な通報先、内部調査通知がない等の追加条件が求められる。。列挙された事情は法定の選択的条件の例示である。全条件の同時成立や、必ず内部を経由する義務とは読まない。\par

\paragraph{再利用する考え方} 公益通報の対象者・通報先別の保護要件と施行時点について、対象・要件・例外・時点を分け、類似用語や無条件の断定を避ける。

\paragraph{共通理論} \hyperref[theory:compliance:whistleblower]{第\ref*{theory:compliance:whistleblower}節}

\par 参考：\cite{caa_handbook_pre2022,caa_report_current,caa_amend2020,caa_schedule2026}。

\par 難易度 3（校準初稿）、優先度未評価。

\subsection{令和2年度 問15}
\label{q:r02-reg-apt-15}

\noindent\textbf{原典：} sources/R02/aptitude.pdf、PDF物理15頁。

\textbf{公式正答：} 1（single）。

\paragraph{問題の要旨} 内部告発の根拠確認、組織手段、同僚の支持、上司への異議、解決策、理解者の6項目で適切な数を求める。

\Needspace{4\baselineskip}\paragraph{確認した条件・記述（独自要約）}\mbox{}\par

\noindent\textbf{方向：} 適切なものの数。6/5/4/3/2。\par

\noindent\textbf{構成：} ア〜カの6項目。\par

\noindent\textbf{限定：} 上司への通知は適当な場合。法的保護要件の全体を問う問題ではない。\par

\noindent\textbf{時点：} 令和2年。2022年以降の改正を出題当時の要件として加えない。\par

\paragraph{問われる核心} 通報の正当な根拠と是正に向けた手続

\paragraph{必要知識} 通報の正当な根拠と是正に向けた手続。ア○：主張の根拠を自ら検証する。あらゆる努力を違法な証拠収集の許可へ拡張しない。イ○：組織内の通報・相談手段を活用する。ウ○：専門職の支持・検討を求め、独断を点検する。エ○：適当な場合という条件付きで直属上司へ異議を伝える。オ○：事実を指摘するだけでなく具体的な是正案を示す。カ○：上司だけで終わらず他の理解者・相談経路を探す。6項目なので1。これは倫理的な行動の検討で、内部通報が外部通報の法的保護の必須前置手続だという意味ではない。

\paragraph{考え方・正答を選ぶ理由（査読待ち初稿）} ア○：主張の根拠を自ら検証する。あらゆる努力を違法な証拠収集の許可へ拡張しない。イ○：組織内の通報・相談手段を活用する。ウ○：専門職の支持・検討を求め、独断を点検する。エ○：適当な場合という条件付きで直属上司へ異議を伝える。オ○：事実を指摘するだけでなく具体的な是正案を示す。カ○：上司だけで終わらず他の理解者・相談経路を探す。6項目なので1。これは倫理的な行動の検討で、内部通報が外部通報の法的保護の必須前置手続だという意味ではない。

\Needspace{4\baselineskip}\paragraph{選択肢の検討（査読待ち初稿）}\mbox{}\par

\noindent\textbf{1：選択対象} 6項目。6項目とも、慎重に根拠を確かめ是正を進める行動である。\par

\noindent\textbf{2：選択対象ではない} 5項目。適切な項目を一つ除外している。\par

\noindent\textbf{3：選択対象ではない} 4項目。適切な項目を二つ除外している。\par

\noindent\textbf{4：選択対象ではない} 3項目。適切な項目を三つ除外している。\par

\noindent\textbf{5：選択対象ではない} 2項目。適切な項目を四つ除外している。\par

\paragraph{再利用する考え方} 対象、条件、時点を先に確認する。通報の正当な根拠と是正に向けた手続を具体的な事例へ適用し、技術的原因・法的要件・倫理的判断のどれを問うか区別する。

\paragraph{共通理論} \hyperref[theory:compliance:whistleblower]{第\ref*{theory:compliance:whistleblower}節}, \hyperref[theory:ethics:disclosure]{第\ref*{theory:ethics:disclosure}節}

\par 参考：\cite{ipej_ethics2011,caa_handbook_pre2022,caa_current_faq}。

\par 難易度 2（校準初稿）、優先度未評価。
